% ---------------------------------------------------------
% 2. SISTEMA DE NUMERACIÓN INTELIGENTE (NUEVO)
% ---------------------------------------------------------
\makeatletter % Permite usar comandos internos de LaTeX con arroba (@)

% 1. Definimos el contador base dependiendo de la clase del documento
\@ifundefined{c@chapter}{
    \newcounter{numDef}[section] % Si es un artículo, reinicia en cada sección
}{
    \newcounter{numDef}[chapter] % Si es un libro, reinicia en cada capítulo
}

% 2. La lógica condicional: Imprime solo los niveles que existen (>0)
\renewcommand{\thenumDef}{%
    \@ifundefined{c@chapter}{%
        % --- LÓGICA PARA ARTICLE ---
        \ifnum\value{subsection}>0 \thesubsection.\arabic{numDef}%
        \else \ifnum\value{section}>0 \thesection.\arabic{numDef}%
        \else \arabic{numDef}\fi\fi
    }{%
        % --- LÓGICA PARA BOOK / REPORT ---
        \ifnum\value{subsection}>0 \thesubsection.\arabic{numDef}%
        \else \ifnum\value{section}>0 \thesection.\arabic{numDef}%
        \else \ifnum\value{chapter}>0 \thechapter.\arabic{numDef}%
        \else \arabic{numDef}\fi\fi\fi
    }%
}
\makeatother % Cerramos el uso de comandos internos

% 3. El comando para insertar la definición sigue igual
\newcommand{\itemdef}[1]{%
    \refstepcounter{numDef}%
    \item[\textbf{\color{primaryColor}Definición \thenumDef \quad (#1):}]%
}

\newtcolorbox{DefinicionesBox}[1][Definiciones]{
    colback=bgColor,            
    colframe=primaryColor,      
    coltitle=accentColor,       
    fonttitle=\bfseries,        
    title=#1,                   
    arc=1mm,                    
    boxrule=1.2pt,              
    left=3mm, right=3mm, top=3mm, bottom=3mm,                 
    drop fuzzy shadow=secondaryColor!30
}